% he-principal-commitments.tex -- Chapter 6 addition (writer: authors).
% Inserted at the end of the section labelled sec:deontic-decision, after
% the reachable-future theorem, its boundary and its exercise pointer.

\subsection{Principal-tagged commitments}\label{sec:he-principal-commitments}
\pdstatus{\Designed}{tagging and the deletion filter are specified, not implemented}

The checker above takes one policy set. In a shared harbor several people
issue commitments into it. Bob forbids migrations during a release freeze;
Alice's direction derives an obligation to migrate. Tag each sentence with the
principal who issued it, authenticated by the card that submitted it. The tag
is metadata. It adds no atom and changes no clause of the conflict definition.
Write $P=P_1\uplus\cdots\uplus P_M$ for the tagged union of the principals'
accepted sentences.

When the tagged set is in conflict, the \pdterm{blame set} is a set of
principals whose commitments together cause the conflict, with no member
redundant; the deletion filter below finds one.
\pdchapref{ls}{The Legible Swarm} uses the blame set
to route escalations%
\ifpdbook
\ (\S\ref{ls:sec:ls-many-authors})%
\fi
.

\begin{pdclaim}{Theorem}{Principal-Tagged Detection and Minimal Blame}
\setlength{\parskip}{4pt plus 1pt}
\label{thm:he-tagged-blame}
Let $P=P_1\uplus\cdots\uplus P_M$ be a policy set in $\mathcal{L}_c$ whose
sentences carry issuing-principal tags.
\begin{enumerate}[leftmargin=1.8em,itemsep=1pt,label=(\alph*)]
\item The checker of Theorem~\ref{thm:deontic-detect} decides whether $P$ is
in conflict within that theorem's bound, and every sentence its witness uses
carries the tag of its issuer.
\item Conflict is upward-closed: if $Q\subseteq P$ is in conflict, so is $P$.
\item Let $P$ be in conflict. The \emph{deletion filter} sets
$B\leftarrow\{1,\dots,M\}$ and then, for $i=1,\dots,M$ in turn, removes $i$
from $B$ whenever $\biguplus_{j\in B\setminus\{i\}}P_j$ is in conflict. It
makes exactly $M$ checker calls and returns a set $B$ such that
$\biguplus_{j\in B}P_j$ is in conflict and no proper subset of $B$ is.
\end{enumerate}
\end{pdclaim}
\begin{proof}
\setlength{\parskip}{4pt plus 1pt}
(a) No clause of the conflict definition reads a tag, so $P$ is an ordinary
member of $\mathcal{L}_c$ and Theorem~\ref{thm:deontic-detect} applies
unchanged. A witness is built from sentences: the facts and rules of a
derivation chain, the rules whose tokens clash, the claims, or the difference
constraints on a cycle. Each sentence keeps its tag.

(b) Adding facts or definite Horn rules can only enlarge the least Horn model,
so a derivation of $\bot$ in $Q$ persists in $P$: clause (i). Bodies are
monotone, so every token fired in $Q$ also fires in $P$, and an overlapping
obligation and prohibition persist: clause (ii). The claims of $Q$ are claims
of $P$: clause (iii). The deadline graph of $Q$ is a subgraph of that of $P$,
so a negative cycle persists: clause (iv). This is the monotonicity of
Theorem~\ref{thm:he-reachable-admission}, applied to sentences rather than
facts.

(c) The union over $B$ is in conflict initially by hypothesis, and $B$ shrinks
only to a set whose union is in conflict, so the final union is in conflict.
Let $i$ remain in the final $B$, and let $B_i$ be the value of $B$ when $i$ was
tested. The union over $B_i\setminus\{i\}$ was conflict-free. The final $B$ is
a subset of $B_i$, so $B\setminus\{i\}\subseteq B_i\setminus\{i\}$, and by
(b) the union over $B\setminus\{i\}$ is conflict-free. Every proper subset of
$B$ lies inside some $B\setminus\{i\}$, so by (b) again it is conflict-free.
The loop makes one checker call per principal.
\end{proof}

A blame set of size one means a principal contradicts their own earlier
commitments, and the conflict goes back to them alone. The set is minimal, not
minimum: different test orders can return different minimal sets
(Exercise~\ref{ex:he-blame-order}). No claim is made here about the complexity
of finding a smallest one. Run over sentences instead of principals, the same
filter returns an irredundant core with $|P|$ calls. A related scheduling tool
extracts unsatisfiable cores from an SMT solver and relaxes the constraints
they name; its authors note that the returned core is not
minimal~\pdcite{kasi2013cassandra}.

The frontier of Theorem~\ref{thm:deontic-frontier} reappears between
principals, in a sharper form.

\begin{pdclaim}{Theorem}{Individually Consistent, Jointly NP-Complete}
\setlength{\parskip}{4pt plus 1pt}
\label{thm:he-joint-frontier}
Extend $\mathcal{L}_c$ with disjunctive obligations under discharge-choice
semantics, as in Theorem~\ref{thm:deontic-frontier}. Deciding whether
$P_1\uplus P_2$ has a conflict-free discharge selection is NP-complete, even
on instances where $P_1$ alone and $P_2$ alone each have one.
\end{pdclaim}
\begin{proof}
\setlength{\parskip}{4pt plus 1pt}
Membership is that of Theorem~\ref{thm:deontic-frontier}: the union is a policy
set of the extended language.

For hardness, take that theorem's reduction from a 3-SAT formula $\varphi$ and
divide its sentences. Give $P_1$ the variable rules
$\mathrm{sel}^{+}_x\to\mathsf{O}(\mathrm{assert}_x)$ and
$\mathrm{sel}^{-}_x\to\mathsf{F}(\mathrm{assert}_x)$, on identical scope and
interval. Give $P_2$ the clause obligations over the selector literals.

$P_1$ alone has no disjunctive obligation and no fact, so its only discharge
selection is empty. No selector holds, no token fires, and $P_1$ has no claim,
deadline or integrity constraint: it is conflict-free. $P_2$ alone has no
prohibition, no rule that a selector can fire, and no claim, deadline or
integrity constraint, so every discharge selection is conflict-free. The union is exactly the reduction's instance, which has a
conflict-free selection if and only if $\varphi$ is satisfiable. The division
takes linear time.
\end{proof}

Each person's commitments make sense on their own; the clash lives only in the
combination, and finding a compatible choice is as hard as satisfiability.
The practical remedy is the one already given for one principal: bounded
search that returns a checkable selection, or a timeout on which the admission
check refuses and attaches what the search found.

Prohibitions behave better than obligations.

\begin{pdclaim}{Theorem}{Prohibitions Compose}
\setlength{\parskip}{4pt plus 1pt}
\label{thm:he-prohibitions-compose}
Let $P\in\mathcal{L}_c$ be conflict-free, and let $\Delta$ be a finite set of
deontic rules of modality $\mathsf{F}$ with monotone bodies. Then
$P\cup\Delta$ is in conflict if and only if some token fired by $\Delta$
clashes with an $\mathsf{O}$ token fired by $P$: same action, overlapping
scope and overlapping interval. In particular, a policy set issued by any
number of principals that contains no $\mathsf{O}$ rule, no integrity
constraint, no exclusive claim and no deadline constraint is conflict-free.
\end{pdclaim}
\begin{proof}
\setlength{\parskip}{4pt plus 1pt}
$\Delta$ adds no fact, Horn rule or integrity constraint, so $P\cup\Delta$ has
the least Horn model of $P$ and clause (i) is unchanged. Claims and the
deadline graph are unchanged, so clauses (iii) and (iv) are unchanged. The
fired tokens are those of $P$ together with $\Delta$'s, all of modality
$\mathsf{F}$. A clause-(ii) clash pairs an $\mathsf{O}$ token with an
$\mathsf{F}$ token; $P$ had none, and every $\mathsf{O}$ token comes from $P$,
so any new clash pairs a token of $\Delta$ with an $\mathsf{O}$ token of $P$.
The converse is immediate. For the second statement: definite Horn rules
without integrity constraints never derive $\bot$; without $\mathsf{O}$ rules
clause (ii) cannot hold; without claims clause (iii) cannot; and an empty
deadline graph has no cycle.
\end{proof}

The admission check for a new prohibition from anyone therefore needs only an
index lookup against the obligations fired on the same action, at the current
facts and at each maximal member of the declared reachable family. A new fact,
Horn rule, obligation, claim or deadline needs the full check, because it can
fire or meet anyone's rules.
\pdscope{The theorem does not say that prohibitions are harmless.}
They can forbid every way to discharge an obligation, which is a clause-(ii)
clash and is caught. They can also forbid so much that nothing useful remains
allowed, a failure of progress that this fragment does not call a conflict.

\begin{pdexample}{One Running Case Through the Section's Theorems}
\setlength{\parskip}{4pt plus 1pt}
This is a constructed example. Bob files a release freeze,
$\mathsf F(\texttt{migrate},\{\texttt{db:prod}\},[0,48])$. Alice files
$\texttt{schema\_approved}\to\mathsf O(\texttt{migrate},\{\texttt{db:prod}\},[24,30])$.
Today \texttt{schema\_approved} is not a fact, so the batch check passes.

\emph{Reachable futures} (Theorem~\ref{thm:he-reachable-admission}). The
declared reachable fact sets are $\varnothing$ and
$\{\texttt{schema\_approved}\}$. The theorem says to check only the largest.
There the obligation fires and overlaps the freeze on
$[\max(0,24),\min(48,30)]=[24,30]$, so the daemon refuses Alice's rule now
rather than when the schema is approved.

\emph{Blame} (the minimal-blame theorem above). Suppose Alice also issued
the fact \texttt{schema\_approved}, so the combined set is in conflict. Run the
deletion filter in the order Bob, then Alice. Without Bob, Alice's set fires an
obligation and meets no prohibition, so it is conflict-free and Bob stays.
Without Alice, only Bob's prohibition remains, which is conflict-free, so Alice
stays. The blame set is $\{\text{Bob},\text{Alice}\}$, found in two checker
calls, and neither principal conflicts alone.

\emph{Either-or} (Theorem~\ref{thm:deontic-frontier} and its two-principal
form above). Let Alice's duty instead be ``migrate, or
ship behind a feature flag''. Choosing the flag removes the clash, since no
prohibition names that action. With many such choices from several
principals, each consistent alone, finding a joint choice is NP-complete in
general, so the daemon searches with a time limit and refuses on a timeout.

\emph{Prohibitions} (Prohibitions Compose, above). Pike adds
$\mathsf F(\texttt{deploy},\{\texttt{prod}\},[0,60])$. The daemon looks up the
obligations fired on \texttt{deploy}, at today's facts and at
$\{\texttt{schema\_approved}\}$, finds none, and accepts the freeze without a
full re-check.
\end{pdexample}

Supervisory control draws a similar line. A specification there may only
restrict what can happen, as a prefix-closed language (one in which every
prefix of an allowed run is also allowed), or may also mark completions that
must stay reachable. Prefix-closed languages are trivially
nonconflicting, while verifying that a synchronous product of an unspecified
number of components is nonconflicting is
PSPACE-complete~\pdcite{komenda2016distributed}; many modular supervision
problems share that complexity~\pdcite{rohloff2005pspace}. The parallel is
structural, not a transfer of either result.

\pdexercisepointer{\ref{ex:he-blame-order}}{\ref{ex:he-new-direction-check}}{\pageref{ex:he-blame-order}}
